\begin{appendices}
\setcounter{section}{2} 

%%%%%%%%==================================================%%%%%
%%%%% MATH PART %%%%%%%
\section{Theoretical foundation of SatuTe}\label{suppsec:maths}
\subsection{Basic properties of the coherence coefficient}
\label{sec:apx}

As inferred from Eq.\,\eqref{eq:branchlikelihod_new}, for pattern $\partial$ the coherence coefficient corresponding to eigenvalue $\lambda_i$ is defined as
\begin{equation} \label{eq:coherence_again}
    C^\partial_i := \langle \frac{\va}{\Prob(\pred)},\bm{h_i}\rangle \langle \frac{\vb}{\Prob(\pblue)}, \bm{h_i}\rangle.
\end{equation}

Let us restate the spectral decomposition of Eq.\,\eqref{eq:spectral_decomposition}. Due to reversibility, the transition matrix $\bm{P}(t) = e^{Qt}$ decomposes as 

\begin{align} \label{eq:spectral_decompositionAppendix}
   \bm{P}(t) = \bm{1}\bm{\pi}^T + \sum_{i=1}^3\bm{v_i} \bm{h_i}^T e^{\lambda_it},
\end{align}
 
where  ${\lambda_0 = 0 > \lambda_1 \geq \lambda_2 \geq \lambda_3}$  are the eigenvalues of the rate matrix and $\bm{h_i}, \bm{v_i}$  are the corresponding left and right eigenvectors. 
Note that the right eigenvalues form an orthonormal basis with respect to $\bm{\pi}$-inner product $\langle \bm{v},\bm{w}\rangle_{\bm{\pi}}= \bm{v}^T\cdot\diag(\bm{\pi})\cdot\bm{w}$ and $\bm{h}_k=\diag(\bm{\pi})\bm{v}_k$ \citepMain{levin2017markov}.


\medskip


To simplify the computation of the expected value of the coherence coefficient, in Prop. \ref{prop:totalSumLikelihood} we consider each scalar product of Eq.\,\eqref{eq:coherence_again} independently.

\begin{prop}\label{prop:totalSumLikelihood} Assuming stationarity, for each pattern $\pred$ observed at the tips of subtree $\mathbb{T}_A$, consider the likelihood vector $\va$ and  the probability ${\Prob(\pred) = \langle \bm{\pi},  \va \rangle }$. The following holds.

\smallskip


\begin{enumerate}[label=\alph*)]
\item If $\bm{1}$ denotes the column vector of $1$'s, then
\begin{equation*} \label{eq:expected1}
\mathbb{E}\Big[\frac{\va}{\Prob(\pred)}\Big] = \bm{1}.
 \end{equation*}
\item If vector $\bm{h_i}$ is orthogonal to $\bm{1}$, then it holds that
 \begin{equation*}\label{eq:expectedscalar0}
 \mathbb{E}\Big[\langle \frac{\va}{\Prob(\pred)}, \bm{h_i}\rangle\Big] = 0.
 \end{equation*}
 \end{enumerate}
\end{prop}
\begin{proof} \hfill


\begin{enumerate}[label=\alph*)]
\item It  is enough to show that 
 \begin{equation} \label{eq:sum1}
      \sum _\pred \va = \bm{1}.
 \end{equation}
Consider one nucleotide $i$. Focusing on the $i$th entry of $\va$, we see that \[\sum_\pred \Prob(\pred  \mid  \text{$i$ at node $A$}) = 1\]
because we are summing over all possible outcomes $\pred$, as desired.
\item We just need to use linearity and item $a)$. Indeed,
\begin{equation}
\mathbb{E}\Big[\langle \frac{\va}{\Prob(\pred)}, \bm{h_i}\rangle\Big] =\langle  \mathbb{E}\Big[\frac{\va}{\Prob(\pred)}\Big], \bm{h_i}\rangle = \langle  \bm{1}, \bm{h_i}\rangle = 0,
\end{equation} 
where in last equality we used the fact that $\bm{h_i}$ and $\bm{1}$ are orthogonal, as it is the case in the spectral decomposition of Eq.\,\eqref{eq:spectral_decompositionAppendix}.
%, as stated in Equation \ref{eq:decomposition}.
\end{enumerate}
\end{proof}

Recall that Eq.\,\eqref{eq:centraleq} implies that, when the length of branch $AB$ grows, the subpatterns $\pred$ and $\pblue$ become independent. Using Prop. \ref{prop:totalSumLikelihood}, we can easily compute the expected value $\mathbb{E}[C^\partial_i]$ under the assumption of subtree independence as follows. 

\begin{prop} \label{prop:CoherenceToZero} In a phylogenetic tree $\mathbb{T}$, consider a branch $AB$ with length $t \to \infty$. Assuming stationarity and reversibility, for any $i \in \{1,2,3\}$ it holds that 
    \[ \mathbb{E}[C^\partial_i] = 0.\]
\end{prop}
\begin{proof}
Since subpatterns $\pred$ and $\pblue$ are sampled independently, also $\langle \va/\Prob(\pred),\bm{h_i}\rangle$ and $\langle \vb/\Prob(\pblue), \bm{h_i} \rangle$ are independent. The expectation is multiplicative under independence, and thus we have
 \begin{align}
    \mathbb{E}[C^\partial_i] &= \mathbb{E}\Big[  \langle \frac{\va}{\Prob(\pred)},\bm{h_i}\rangle \langle \frac{\vb}{\Prob(\pblue)}, \bm{h_i} \rangle \Big] =  \\
    &=  \mathbb{E}\Big[  \langle \frac{\va}{\Prob(\pred)},\bm{h_i}\rangle \Big] \mathbb{E}\Big[ \langle \frac{\vb}{\Prob(\pblue)}, \bm{h_i} \rangle \Big] = 0,
\end{align}
where we used Prop. \ref{prop:totalSumLikelihood}b.

\end{proof}

As described in Methods \ref{msec:estimates}, for any $i,j\in \{ 1,2,3\}$ we define the covariance
\begin{equation}\sigma^2_{A,i,j}:=\mathbb{E}\Big[\langle \frac{\va}{\Prob(\pred)},\bm{h_i}\rangle  \langle \frac{\va}{\Prob(\pred)},\bm{h_j}\rangle \Big] 
%= \cov [\langle \frac{\va}{\Prob(\pred)},\bm{h_i}\rangle, \langle \frac{\va}{\Prob(\pred)},\bm{h_j}\rangle]
.\end{equation}

This is indeed the covariance between the two scalar products due to Prop. \ref{prop:totalSumLikelihood}b. As claimed in Methods \ref{msec:estimates}, in the following proposition we prove that $\sigma^2_{A,i,j}$ greatly simplifies if node $A$ is external. 

\begin{prop} \label{prop:memorySequence}
    Assuming stationarity and reversibility, if node $A$ is external (or equivalently, if subtree $\mathbb{T}_A$ is composed by a single sequence), then \[
    \sigma^2_{A,i,j} =  \begin{cases}
    1,  & \text{if $i=j$}.\\
    0,&  \text{if $i\neq j$}.
    \end{cases}
    \]
\end{prop}
\begin{proof}
   In the eigendecomposition of Eq.\,\eqref{eq:spectral_decompositionAppendix}, we know that $\diag(\bm{\pi}) \bm{v_i} = \bm{h_i}$, while $\langle \bm{v_i} ,\bm{h_j} \rangle = 1$ and $\langle \bm{v_i} ,\bm{h_j} \rangle = 0$ if $i \neq j$. The patterns of a single sequence are the possible letters $x$ it is composed of.  Moreover, due to stationarity, we know that pattern $\pred = x$ satisfies $\Prob(\pred) = \pi_x$, where $x$ is a nucleotide. Note also that, if $\pred = x$,  then $\bm{L}(\pred)$ is the vector of $0$'s, except $1$ at the $x$th entry. All in all, we have
\begin{equation}
\sigma^2_{A,i,i} = \sum_x \Prob(x) (\frac{1}{\pi_x} h_i^x )^2 = \sum_x \pi_x (\frac{1}{\pi_x} h_i^x )^2 = \sum_x v_i^x h_i^x = \langle \bm{h_i}, \bm{v_i} \rangle= 1.
\end{equation}
Similarly, if $i \neq j$ we have 
\begin{equation}
  \sigma^2_{A,i,j}   = \sum_x \Prob(x) (\frac{1}{\pi_x} h_i^x )(\frac{1}{\pi_x} h_j^x ) = \sum_x \pi_x (\frac{1}{\pi_x} h_i^x )(\frac{1}{\pi_x} h_j^x ) = \langle \bm{h_i}, \bm{v_j} \rangle = 0.
\end{equation}
\end{proof}


\subsection{Testing for Saturation for Any Multiplicity of \texorpdfstring{$\lambda_1$}{l1}}
\label{sec:multiplicityGreaterThan1}

Any linear combination of the coherence coefficients them can be used to test for saturation. However, we want the most powerful test among the candidates. As we prove in Appendix \ref{subsec:approximation}, optimal power is achieved asymptotically by the test that uses the coherence coefficient introduced in Eq.\,\eqref{eq:estimator_explicit}, namely
\begin{equation}\label{eq:dominant_coherence}
\hat{C}_1:= \sum_{k \in \{ 1, \cdots , D \}}\hat{C}_{1, k},
\end{equation} 
where eigenvalue $\lambda_1$ has multiplicity $D$ and $\lambda_1 = \cdots = \lambda_D$. Under the null hypothesis of subtree independence, the expectation and variance of $\hat{C}_1$ can be computed as described in the following proposition.

\medskip
 
\begin{prop} \label{prop:variance_coherence_D}  In a phylogenetic tree $\mathbb{T}$ consider a branch $AB$ with length $t^*$ and assume stationarity and reversibility.  If $t^* \to \infty$, then the coherence coefficient $\hat{C}_1$ satisfies
\begin{align}
 \mathbb{E}[\hat{C}_1] &= 0,\\
 \mathrm{Var}[\hat{C}_1] &= \frac{1}{n}\sum_{i, j \in \{1,\cdots, D\}} 
   \sigma^2_{A,i,j}   \sigma^2_{B,i,j}.
\end{align}
In particular, if node $A$ is external, then 
\begin{equation} \label{eq:externalMultiple}
 \mathrm{Var}[\hat{C}_1]= \frac{1}{n}\sum_{i \in \{1,\cdots, D\}} 
      \sigma^2_{B, i, i},
\end{equation}
and if both nodes $A$ and $B$ are external, then 
\begin{equation}
 \mathrm{Var}[\hat{C}_1] = \frac{D}{n}.
 \end{equation}
\end{prop}
\begin{proof} Since the expectation of a sum is the sum of expectations, $\mathbb{E}[C_{1,i}^\partial] = 0$ (Prop. \ref{prop:CoherenceToZero}) implies that  $\mathbb{E}[\hat{C}_1] = 0$. 
 
 \medskip
 
 Regarding the variance, we can sum up the variance of each independent site, giving $\mathrm{Var}[\hat{C}_1] = 1/n  \mathrm{Var}[\sum_{i \in \{1,\cdots,D\}}C^\partial_{1,i}] $. Since $\mathbb{E}[C^\partial_{1,i}] = 0$ (Prop. \ref{prop:CoherenceToZero}), we have 
 \begin{equation} \label{eq:VarianceExpected}
     \mathrm{Var}[\sum_{i \in \{1,\cdots, D\}}C^\partial_{1,i}] =\mathbb{E}[(\sum_{i \in \{1,\cdots, D\}}C^\partial_{1,i})^2 ].
 \end{equation}
 After expanding the squared sum $(\sum_{i \in \{1,\cdots, D\}}C^\partial_{1,i})^2$, since the expectation is multiplicative under independence, a summand 
 $C^\partial_{1,i} C^\partial_{1,j}$ satisfies
  \begin{align}
 &\mathbb{E}[C^\partial_{1,i} C^\partial_{1,j}] = \\
 =& \mathbb{E}\Big[\langle \frac{\va}{\Prob(\pred)},\bm{h_i}\rangle \langle \frac{\vb}{\Prob(\pblue)}, \bm{h_i} \rangle \langle \frac{\va}{\Prob(\pred)},\bm{h_j}\rangle \langle \frac{\vb}{\Prob(\pblue)}, \bm{h_j} \rangle\Big]= \nonumber\\
  =& \mathbb{E}\Big[\langle \frac{\va}{\Prob(\pred)},\bm{h_i}\rangle \langle \frac{\va}{\Prob(\pred)},\bm{h_j}\rangle \Big]  \mathbb{E}\Big[\langle \frac{\vb}{\Prob(\pblue)}, \bm{h_i} \rangle \langle \frac{\vb}{\Prob(\pblue)}, \bm{h_j} \rangle\Big]= \nonumber\\
  =& \sigma^2_{A,i,j}   \sigma^2_{B,i,j}.\nonumber
 \end{align}
 Summing for all $i,j \in \{1,\cdots, D\}$ we get the desired equality.  The particular case when nodes $A$ and/or $B$ are external follows from Prop. \ref{prop:memorySequence}.
 
\end{proof}


Prop. \ref{prop:variance_coherence_D} leads to a one-sided $z$-test based on the coherence coefficient $\hat{C}_1$. The null hypothesis of independence is rejected with significance $\alpha$ if 
\begin{equation} \label{eq:threshold_D}
     \hat{C}_1 >  z_\alpha \sqrt{
     {\mathrm{\widehat{Var}}}[\hat{C}_1]
     } = z_\alpha \frac{\sqrt{\sum_{i, j \in \{1,\cdots, D\}} 
   \hat{\sigma}^2_{A,i,j}   \hat{\sigma}^2_{B,i,j}.
}}{\sqrt{n}},
\end{equation} where $z_\alpha$ satisfies $\Prob(Z > z_\alpha) = \alpha$, while $\hat{\sigma}^2_{A,i,j}$ is the sample estimate of $\sigma^2_{A,i,j}$ introduced in Eq.\,\eqref{variance_estimate}. If node $A$ is external, then Eq. \ref{eq:threshold_D} is simplified due to Prop. \ref{prop:memorySequence} as 
\begin{equation} \label{eq:threshold_D_external}
     \hat{C}_1 > z_\alpha \frac{\sqrt{\sum_{i\in \{1,\cdots, D\}}  \hat{\sigma}^2_{B,i,i}
}}{\sqrt{n}}.
\end{equation}
\bigskip

\subsection{The Asymptotic Equivalence between the MLE and the Dominant Coherence
}
\label{sec:asym_phylo_tree}

There is a close relationship between the coherence coefficient and the log-likelihood of a branch. To describe this relationship, consider branch $AB$ of tree $\mathbb{T}$. Given an alignment where pattern $\partial$ is observed $n_\partial$ times, the log-likelihood $f(t)$ of branch $AB$ having length $t$ is
\begin{equation}\label{eq:definitionloglike}
f(t) = 
\sum_\partial n_\partial \log{(\Prob(\partial \mid t))}.
\end{equation}
If we define the log-likelihood of independence as 
 \begin{equation} \label{eq:DIFLoglikelihood}
 {f_\infty := \sum_\partial n_\partial \Big(\log\Prob(\pred)+  \log\Prob(\pblue)\Big)}, 
 \end{equation}
 then Eq.\,\eqref{eq:definitionloglike} can be rewritten as 
\begin{equation}\label{eq:definitionloglike2}
f(t) -f_\infty= 
\sum_\partial n_\partial \log{(\frac{\Prob(\partial \mid t)}{\Prob(\pred)\Prob(\pblue)})}.
\end{equation}
    Using Eq.\,\eqref{eq:centraleq}, last equation becomes
 \begin{align} 
     \label{eq:efficient_likelihood_2}
     f(t)-f_\infty &=  \sum_{\partial} n_\partial \log \Big( 1+  \sum_{i \in \{1,2,3\}} e^{\lambda_it}  C_i^\partial\Big).
 \end{align} 
 
In Prop. \ref{prop:limit_likelihood_tree},  we formalise the relationship between the coherence coefficient and the log-likelihood $f(t)$. 
 
  \begin{prop} \label{prop:limit_likelihood_tree}
Assuming stationarity and reversibility, if rate matrix $Q$ has eigenvalues $0 > \lambda_1 \geq \lambda_2 \geq \lambda_3$ and $\lambda_1$ has multiplicity $D$, then define the coherence coefficient \[\hat{C}_1 := \sum_{i \in \{1,\cdots, D\}} \hat{C}_{1,i}.\]
If $\hat{C}_1 \neq 0,$ then \[f(t) - f_\infty \sim \hat{C}_1 n e^{\lambda_1 t}.\]
\end{prop}
\begin{proof} \hfill
 \begin{sloppypar}
 In Equation \ref{eq:definitionloglike2} it is clear that $f(t) - f_\infty$ tends to zero as $t \to \infty$, since $\log(1) = 0$. Now, using L'H\^{o}pital's rule, it is enough to show that the derivative $f'(t)$ satisfies ${f'(t) \sim 
     \hat{C}_1 n \lambda_1  e^{\lambda_1 t}}$. In the expression
 \end{sloppypar}
 \begin{equation*}
     [\log \Big(1 +  \sum_{i \in \{ 1,2,3\}} e^{\lambda_i t} C_i^\partial \Big)]' = \frac{  \sum_{i \in \{ 1,2,3\}} \lambda_i e^{\lambda_it} C_i^\partial }{1 + \sum_{i \in \{ 1,2,3\}} e^{\lambda_it} C_i^\partial},
 \end{equation*}
 the denominator is asymptotically equivalent to $1$. On the other hand, if ${\sum_{i \in \{1,\cdots, D\}}C_i^\partial \neq 0,}$ then the numerator is equivalent to ${\lambda_1e^{\lambda_1t}\sum_{i \in \{1,\cdots, D\}}C_i^\partial}$, and $o(e^{\lambda_1t})$ otherwise. Due to the assumption $\hat{C}_1 \neq 0$, we can sum up the equivalencies, giving
\begin{equation}
     f'(t) \sim \sum_{\partial} n_\partial \Big(
     \lambda_1 e^{\lambda_1t}
     \sum_{i \in \{1,\cdots, D\}}
     C_{1,i}^\partial 
     \Big) = 
     n \lambda_1  e^{\lambda_1 t} \sum_{i \in \{1,\cdots, D\}} 
     \hat{C}_{1,i}= n \lambda_1  e^{\lambda_1 t} \hat{C}_1,
\end{equation}
as desired.

\end{proof}

\iffalse
Notably, Prop. \ref{prop:limit_likelihood_tree} implies that whether $f(t)$ increases or decreases for large $t$ uniquely depends on the sign of $\hat{\delta}$, as stated in the following corollary. 

\begin{cor} \label{prop:decission_tree}
\begin{sloppypar}
Given an alignment and assuming a stationary and reversible process on a tree, the log-likelihood $f(t)$ of branch $AB$ having length $t$ satisfies the following.

\end{sloppypar}
\begin{itemize}
    \item[a)] If $\hat{\delta}$ is positive, then $f(t)$ has local minimum $f_\infty$ when $t \to \infty$.
    \item[b)] If $\hat{\delta}$ is negative, then $f(t)$ has local maximum $f_\infty$ when $t \to \infty$.
\end{itemize}
\end{cor}
Corollary \ref{prop:decission_tree} is an additional motivation (apart from the approximation of the likelihood-ratio test in Subappendix \ref{subsec:approximation}) to perform a one-sided test for saturation in Equations \ref{eq:threshold} and \ref{eq:threshold_D}. Indeed, it would be contradictory to reject saturation because  $\hat{\delta} \ll 0$, while the MLE could be $\hat{t} = \infty$ due to Cor. \ref{prop:decission_tree}b. 

\fi

%Leaving aside our test for saturation, we recommend using Corollary \ref{prop:decission_tree} to predict the possible numerical divergence of the Newton method. If the log-likelihood $f(t)$ has a unique maximum for $t \in [0, \infty]$, Corollary \ref{prop:decission_tree} characterizes the convergence of the Newton method  in the search for the absolute maximum of $f(t)$. %as explained in Appendix \ref{sec:multiple_maxima}. For practitioners, 
 %However, in general Corollary \ref{prop:decission_tree} cannot be used to guarantee that an MLE $\hat{t} < \infty$ is indeed the absolute maximum point, since the log-likelihood $f(t)$ can have multiple maxima.
 
 \subsection{Optimal power of the test for saturation}
\label{subsec:approximation}
We know that the likelihood-ratio test is the most powerful $\alpha$-level test, as proved by \citeSupp{Neyman1933}. 
%In Appendix  \ref{sec:asym_phylo_tree}, we have described the the log-likelihood $f(t)$ of branch $AB$ having length $t$, which has limit $f_\infty$ as $t \to \infty$ (Equations \ref{eq:definitionloglike}, \ref{eq:DIFLoglikelihood}, and \ref{eq:definitionloglike2}). 
Given the log-likelihood $f(t)$ of branch $AB$ having length $t$, if the MLE is $\hat{t} < \infty$, the likelihood-ratio test compares the null hypothesis $t^* \to \infty$ versus the alternative hypothesis $t^* = \hat{t}$  using statistic $f(\hat{t})-f_\infty$, or more explicitly
\begin{equation}
    \text{"Reject $t^* \to \infty$ if $f(\hat{t})-f_\infty > c$",}
\end{equation}
 where $c \in \mathbb{R}_{>0}$ is chosen so that $\Prob(f(\hat{t})-f_\infty > c  \mid  t^* \to \infty) = \alpha$.


\medskip
Assuming that $\hat{t}$ is large enough,
%from the singularity of $f(t)$, 
we can use the approximation ${f(\hat{t}) \approx  f_\infty + \hat{C}_1 n e^{\lambda_1 \hat{t}}}$ of Prop. \ref{prop:limit_likelihood_tree}. This gives the test 
\begin{equation}
    \text{"Reject $t^* \to \infty$ if $ \hat{C}_1  > e^{-\lambda_1 \hat{t}}c/n$",}
\end{equation}
 where $c \in \mathbb{R}_{>0}$ is chosen so that $\Prob( \hat{C}_1 > e^{-\lambda_1 \hat{t}}c/n \mid  t^* \to \infty) = \alpha$. Setting $c' := e^{-\lambda_1 \hat{t}}c/n$, it is clear that the one-sided test for saturation using the coherence coefficient $\hat{C}_1$ achieves optimal power for long branches. 


\bibliographystyleSupp{apalike}
\bibliographySupp{biblio.bib}


\end{appendices}
